% ============================================================================
% COVER LETTER TEMPLATE (for journal submission: Involve / Experimental Mathematics / AMM)
% ============================================================================
% Dear Editor,
%
% Please consider the attached manuscript, "On a Remarkable Property of the
% Cuboctahedral Rotator, Euler's Identity, and Closed-Form Evaluation of Fresnel,
% Jacobi, Bessel, and Dirichlet Integrals in Lean 4", for publication in [Journal].
%
% Extending the piecewise-linear trigonometric framework studied by Thompson
% (arXiv:1205.3814), where derivatives lost the classical sine-cosine coupling
% at diamond corners and integral calculus remained open, we show that on the
% 12-vertex cuboctahedron O_2(d) (d = 3) and the associated discrete phase ring
% Z_{N_ring} (N_ring = 4*Q_square = 512, Q_square = 128), the piecewise-linear
% projection pair E_d(phi) = (C_d(phi), S_d(phi)) forms an exact discrete Euler
% exponential generating the five constants {0, 1, i, pi, e}:
%   1. A quarter-period shift +Q_square(d) = +128 acts across all 512 points as
%      exact multiplication by the imaginary unit i;
%   2. A half-period shift +256 satisfies Euler's identity e^(i*pi) + 1 = 0
%      identically across the ring;
%   3. The one-step difference Wronskian W_d(phi) = C_d(phi)*S_d(phi+1) -
%      S_d(phi)*C_d(phi+1) = <i*E_d(phi), Delta E_d(phi)> is identically equal
%      to the constant Q_square(d) = 128 at all 512 points, including the four
%      corner vertices.
%
% Using this constant Wronskian, the two-tick Verlet recurrence (unifying
% Liouville's circular basis at zero acceleration with the quadratic-phase
% Fresnel/Airy basis at constant acceleration), and the lemniscate Jacobi
% elliptic triad (sn_d, cn_d, dn_d) at the cuboctahedron modulus k^2 = 1/2,
% we evaluate in exact closed integer form the discrete Fresnel diffraction
% integrals (C_Fresnel = S_Fresnel = 2048), complete and incomplete elliptic
% integrals of the second kind, the Dirichlet sine integral, the Parseval
% energy integral (2,796,288), and the zero-order Bessel integral (with exact
% root at k = 4). Finally, we show that the continuum circle and pi emerge via
% rectification of the 512-gon Weierstrass rotator (with an intrinsic 6.3 ppm
% chord-arc defect on the fundamental ring), while projecting the 12 vertices
% of the cuboctahedron onto any coordinate axis yields the Archimedes equizonal
% partition (4, 4, 4), which together with the central node of weight 12 derives
% Simpson's parabolic weights (4, 16, 4) = 4*(1, 4, 1) and the exact 3D continuum
% Riemann integral from pure 3D geometry.
%
% Every definition, lemma, and theorem is formally verified in the Lean 4 theorem
% prover in a self-contained 495-line module (`RotatorCalculusPaper.lean`, zero
% `sorry`, zero `Classical.choice`) available at:
% https://github.com/KarataevDmitry/cuboctahedral-rotator-calculus
%
% The manuscript is original, has not been published previously, and is not under
% consideration elsewhere.
%
% Sincerely,
% Dmitrii A. Karataev
% ============================================================================

\documentclass[11pt,a4paper]{article}
\usepackage[margin=2.4cm]{geometry}
\usepackage{amsmath,amssymb,amsthm,mathtools}
\usepackage{booktabs}
\usepackage{enumitem}
\usepackage{xcolor}
\usepackage{hyperref}
\usepackage{fontspec}
\setmainfont{Times New Roman}
\setmonofont{Consolas}[Scale=0.90]

\hypersetup{
  colorlinks=true,
  linkcolor=blue!70!black,
  citecolor=blue!70!black,
  urlcolor=blue!70!black,
  pdftitle={On a Remarkable Property of the Cuboctahedral Rotator, Euler's Identity, and Closed-Form Evaluation of Fresnel, Jacobi, Bessel, and Dirichlet Integrals in Lean 4},
  pdfauthor={Dmitrii A. Karataev}
}

\theoremstyle{plain}
\newtheorem{theorem}{Theorem}[section]
\newtheorem{lemma}[theorem]{Lemma}
\newtheorem{proposition}[theorem]{Proposition}
\newtheorem{corollary}[theorem]{Corollary}
\theoremstyle{definition}
\newtheorem{definition}[theorem]{Definition}
\newtheorem{remark}[theorem]{Remark}

\title{\textbf{On a Remarkable Property of the Cuboctahedral Rotator, Euler's Identity, and Closed-Form Evaluation of Fresnel, Jacobi, Bessel, and Dirichlet Integrals in Lean~4}}
\author{Dmitrii A.~Karataev}
\date{}

\begin{document}
\maketitle

\begin{abstract}
We investigate discrete trigonometric and integral calculus on the three-dimensional cuboctahedral shell $O_2(d)\subset\{-1,0,+1\}^d$ ($d=3$, $|V_{O_2}|=12$ vertices, $|\square|=6$ square faces, $|\triangle|=8$ triangular faces) and the associated finite phase ring $\mathbb{Z}_{N_{\mathrm{ring}}}$ ($N_{\mathrm{ring}} = 4\,Q_\square = 512$, where $Q_\square = 2^{|\triangle|-1} = 128$ is the quarter-period). Extending the piecewise-linear calculus of Thompson (\texttt{arXiv:1205.3814}), where piecewise derivatives lost the classical sine--cosine coupling at diamond corners and integral calculus remained open, we prove that the branchless projection pair $\mathcal{E}_d(\varphi) = (\mathcal{C}_d(\varphi),\mathcal{S}_d(\varphi))$ forms an exact discrete Euler exponential on the $512$-gon generating the five fundamental constants $\{0,1,i,\pi,e\}$: a quarter-period shift $+Q_\square(d)=+128$ acts across all $512$ ring points as exact multiplication by the imaginary unit~$i$, a half-period shift $+\Phi_{\mathrm{half}}(d)=+256$ satisfies Euler's identity $\mathcal{E}_d(\varphi+\Phi_{\mathrm{half}}(d))+\mathcal{E}_d(\varphi)=(0,0)$ ($e^{i\pi}+1=0$), and the one-step difference Wronskian $\mathcal{W}_d(\varphi)=\mathcal{C}_d(\varphi)\mathcal{S}_d(\varphi+1)-\mathcal{S}_d(\varphi)\mathcal{C}_d(\varphi+1)=\langle \mathcal{E}_d(\varphi+Q_\square(d)),\Delta\mathcal{E}_d(\varphi)\rangle$ is identically equal to the constant $Q_\square(d)=128$ at all $512$ points of the ring, including the four corner vertices.

We further prove that the two-tick symplectic Verlet difference recurrence $\varphi_{t+1}-2\varphi_t+\varphi_{t-1}=\Phi$ unifies Liouville's elementary circular basis ($\Phi=0$) with the non-elementary quadratic-phase Fresnel--Airy basis ($\Phi\ne 0$), while two-ring coupling at the intrinsic cuboctahedron modulus $k^2 = 1/\|e\|_2^2 = 1/(d-1) = 1/2$ generates Gauss's lemniscate Jacobi elliptic triad $(\mathrm{sn}_d,\mathrm{cn}_d,\mathrm{dn}_d)$. This structure yields exact closed-form integer evaluations of the discrete Fresnel diffraction integrals ($\mathcal{C}_{\mathrm{Fresnel}}=\mathcal{S}_{\mathrm{Fresnel}}=2048$ via half-period cancellation of $480$ oscillating quadratic residues and survival of $32$ stationary-phase points on the bisector $\varphi^2\equiv 64\pmod{512}$), complete and incomplete elliptic integrals of the second kind, the Dirichlet sine integral ($\mathrm{Si}_3(128)=128$), the Parseval energy integral ($2,796,288$), and the zero-order Bessel integral ($\mathcal{J}_{0,3}(0)=65,536$, $\mathcal{J}_{0,3}(1)=32,768$, and exact root $\mathcal{J}_{0,3}(4)=0$). Finally, we show that the continuum circle and $\pi$ emerge as the rectification limit of the $512$-gon Weierstrass rotator (with an intrinsic finite perimeter defect of $\pi^2/(6\cdot 512^2)\approx 6.3\;\text{ppm}$ on the fundamental ring), and that projecting the $12$ vertices of the cuboctahedron $O_2(3)$ onto any coordinate axis yields the Archimedes equizonal partition $(4,4,4)$, which together with the central node of weight $12$ derives the parabolic quadrature weights $(4,16,4)=4\times(1,4,1)$ and the exact 3D continuum Riemann integral $\iiint_{[-1,1]^3} F\,\mathrm{d}V$ from pure 3D geometry. All definitions and theorems are formally verified in the \textbf{Lean~4} theorem prover in a self-contained module without non-constructive axioms (\texttt{0 sorry}, \texttt{0 Classical.choice}).

\medskip\noindent
\textbf{Keywords:} piecewise-linear trigonometry, Euler's identity, difference Wronskian, Weierstrass substitution, cuboctahedron, Fresnel integrals, lemniscate Jacobi elliptic functions, formal verification, Lean~4.

\smallskip\noindent
\textbf{MSC 2020:} 26A06, 33E05, 39A12, 65D30, 68V20.
\end{abstract}

\section{Introduction}
\label{sec:intro}

Piecewise-linear trigonometry in the $L^1$ metric (``taxicab trigonometry'' on the diamond $|x|+|y|=R$) was developed by Krause~\cite{Krause1986} and Thompson~\cite{Thompson2011,Thompson2012}. In~\cite{Thompson2012} (\texttt{arXiv:1205.3814}), Thompson investigated the piecewise derivatives of the $L^1$ sine, cosine, tangent, and secant functions. However, three structural obstacles remained unresolved in~\cite{Thompson2012}:
\begin{enumerate}[leftmargin=*,noitemsep]
\item the piecewise-constant derivatives of the $L^1$ sine and cosine are undefined at the four corner extrema and lose the classical Euclidean differential coupling $(\sin\theta)'=\cos\theta$, $(\cos\theta)'=-\sin\theta$;
\item representing the piecewise branches across quadrants required complex powers $i^{k-2}$ branched by the parity of the quadrant index $k$;
\item integral calculus, the lift to the Euclidean ($L^2$) circle, and the evaluation of classical non-elementary integrals (Fresnel, Jacobi elliptic, Bessel, Dirichlet) were left open.
\end{enumerate}

In this paper, we resolve these obstacles by formulating the rotator over the three-dimensional Archimedean cuboctahedral shell $O_2(d) = \{e=(e_x,e_y,e_z)\in\{-1,0,+1\}^3 \mid e_x^2+e_y^2+e_z^2 = d-1\}$ ($d=3$) and the associated modular integer phase ring $\mathbb{Z}_{N_{\mathrm{ring}}(d)}$:
\[
  |{\triangle}(d)| = 2^d = 8,\qquad
  |\square(d)| = 2d = 6,\qquad
  |V_{O_2}(d)| = 2d(d-1) = 12,
\]
\[
  d_{\mathrm{in}}(\square) = 4,\qquad
  d_{\mathrm{out}}(E)(d) = 2^d - d = 5,\qquad
  \|e\|_2^2 = d - 1 = 2,
\]
\[
  Q_\square(d) = 2^{|{\triangle}(d)|-1} = 128,\qquad
  \Phi_{\mathrm{half}}(d) = 2\,Q_\square(d) = 256,\qquad
  N_{\mathrm{ring}}(d) = 4\,Q_\square(d) = 512,\qquad
  K_P(d) = \frac{N_{\mathrm{ring}}(d)}{|{\triangle}(d)|} = 64.
\]
All definitions and theorems below are formally verified in the self-contained \textbf{Lean~4}~\cite{Moura2021} module \texttt{RotatorCalculusPaper.lean}.

\section{Discrete Euler Exponential, Euler's Identity \texorpdfstring{$e^{i\pi}+1=0$}{e\^{}(i pi)+1=0}, Constant Wronskian, and the Weierstrass 512-gon Lift}
\label{sec:wronskian-euler}

For any $x\in\mathbb{Z}$, define the principal symmetric branch wrapper into $(-\Phi_{\mathrm{half}}(d),\Phi_{\mathrm{half}}(d)]$:
\begin{equation}
  \label{eq:wrap-ring}
  \mathrm{wrap}_d(x) \;:=\; \Phi_{\mathrm{half}}(d) - \bigl((\Phi_{\mathrm{half}}(d) - x) \bmod N_{\mathrm{ring}}(d)\bigr),
\end{equation}
and the branchless piecewise-linear cosine and sine projections:
\begin{equation}
  \label{eq:proj-cs}
  \mathcal{C}_d(\varphi) \;:=\; Q_\square(d) - |\mathrm{wrap}_d(\varphi)|,\qquad
  \mathcal{S}_d(\varphi) \;:=\; \mathcal{C}_d\bigl(\varphi - Q_\square(d)\bigr).
\end{equation}

\begin{definition}[Discrete Euler exponential and one-step difference Wronskian]
\label{def:euler-exp-wronskian}
We define the \emph{discrete Euler exponential} on $\mathbb{Z}_{N_{\mathrm{ring}}(d)}$ by the pair
\[
  \mathcal{E}_d(\varphi) \;:=\; \bigl(\mathcal{C}_d(\varphi),\,\mathcal{S}_d(\varphi)\bigr) \;\in\; \mathbb{Z}^2,
\]
and the \emph{one-step difference Wronskian} by
\begin{equation}
  \label{eq:rotator-wronskian}
  \mathcal{W}_d(\varphi) \;:=\; \mathcal{C}_d(\varphi)\,\mathcal{S}_d(\varphi+1) - \mathcal{S}_d(\varphi)\,\mathcal{C}_d(\varphi+1)
  \;=\; \mathcal{C}_d(\varphi)\,\Delta\mathcal{S}_d(\varphi) - \mathcal{S}_d(\varphi)\,\Delta\mathcal{C}_d(\varphi),
\end{equation}
where $\Delta f(\varphi) = f(\varphi+1)-f(\varphi)$.
\end{definition}

To each $\varphi\in\mathbb{Z}_{N_{\mathrm{ring}}(d)}$ we associate the rational Brahmagupta--Weierstrass lift onto the Euclidean circle $X^2+Y^2=R^2$:
\begin{equation}
  \label{eq:weierstrass-lift}
  X_d(\varphi) := \mathcal{C}_d(\varphi)^2 - \mathcal{S}_d(\varphi)^2,\qquad
  Y_d(\varphi) := \|e\|_2^2\,\mathcal{C}_d(\varphi)\,\mathcal{S}_d(\varphi),\qquad
  R_d(\varphi) := \mathcal{C}_d(\varphi)^2 + \mathcal{S}_d(\varphi)^2.
\end{equation}

\begin{theorem}[Discrete Euler's Identity, Universal Constant Wronskian, Weierstrass Lift, and Fundamental Theorem of Discrete Calculus]
\label{thm:main-1}
For $d=3$, the following exact identities hold across all $N_{\mathrm{ring}}(d)=512$ points $\varphi\in\mathbb{Z}_{N_{\mathrm{ring}}(d)}$ (\texttt{thm1\_constant\_wronskian\_weierstrass\_lift\_and\_newton\_leibniz}):
\begin{enumerate}[label=\textbf{(\roman*)},leftmargin=*]
\item \textbf{Generation of $\{0,1,i,\pi,e\}$ and Euler's identity $e^{i\pi}+1=0$:}
  \begin{itemize}[noitemsep]
  \item the $L^1$ norm is constant: $|\mathcal{C}_d(\varphi)| + |\mathcal{S}_d(\varphi)| \equiv Q_\square(d)$, with $\mathcal{E}_d(0) = (Q_\square(d),0)$ representing unity~$1$;
  \item a quarter-period shift $+Q_\square(d)$ acts across all $512$ points as exact multiplication by the imaginary unit~$i$:
    \[
      \mathcal{E}_d\bigl(\varphi + Q_\square(d)\bigr) \;=\; \bigl(-\mathcal{S}_d(\varphi),\; +\mathcal{C}_d(\varphi)\bigr) \;=\; i\cdot\mathcal{E}_d(\varphi);
    \]
  \item a half-period shift $+2Q_\square(d)=+\Phi_{\mathrm{half}}(d)$ ($\pi$) satisfies Euler's identity $e^{i(\varphi+\pi)}+e^{i\varphi}=0$ identically across the ring:
    \[
      \mathcal{E}_d\bigl(\varphi + \Phi_{\mathrm{half}}(d)\bigr) + \mathcal{E}_d(\varphi) \;=\; (0,0),
      \qquad
      \mathcal{E}_d\bigl(\Phi_{\mathrm{half}}(d)\bigr) + \mathcal{E}_d(0) \;=\; (0,0);
    \]
  \item the inner product of $i\mathcal{E}_d(\varphi)=\mathcal{E}_d(\varphi+Q_\square(d))$ with the difference derivative $\Delta\mathcal{E}_d(\varphi)$ coincides with the Wronskian $\mathcal{W}_d(\varphi)$ and is \textbf{identically equal to the constant $Q_\square(d)$ at all $512$ points of the ring (including the four corner vertices)}:
    \begin{equation}
      \label{eq:constant-wronskian-all-512}
      \bigl\langle \mathcal{E}_d\bigl(\varphi+Q_\square(d)\bigr),\,\Delta\mathcal{E}_d(\varphi)\bigr\rangle
      \;=\; \mathcal{W}_d(\varphi) \;\equiv\; Q_\square(d) \quad (=128),\qquad
      \sum_{\varphi=0}^{N_{\mathrm{ring}}(d)-1}\mathcal{W}_d(\varphi) = N_{\mathrm{ring}}(d)\,Q_\square(d).
    \end{equation}
  \end{itemize}
\item \textbf{Root-free Euclidean Weierstrass $512$-gon and arc-step differential:}
  at all $N_{\mathrm{ring}}(d)$ points, $X_d(\varphi)^2 + Y_d(\varphi)^2 = R_d(\varphi)^2$, and the one-step Euclidean cross and dot products are governed by $\mathcal{I}_d(\varphi) := \mathcal{C}_d(\varphi)\mathcal{C}_d(\varphi+1) + \mathcal{S}_d(\varphi)\mathcal{S}_d(\varphi+1)$:
  \begin{equation}
    \label{eq:brahmagupta-step}
    X_d(\varphi)Y_d(\varphi+1) - Y_d(\varphi)X_d(\varphi+1) \;=\; \|e\|_2^2\,Q_\square(d)\,\mathcal{I}_d(\varphi),\qquad
    X_d(\varphi)X_d(\varphi+1) + Y_d(\varphi)Y_d(\varphi+1) \;=\; \mathcal{I}_d(\varphi)^2 - Q_\square(d)^2.
  \end{equation}
\item \textbf{Fundamental Theorem of Discrete Calculus (Newton--Leibniz):}
  the piecewise-quadratic antiderivatives $\widehat{\Phi}_{\mathcal{C}}(d,\Phi) := \sum_{\varphi=0}^{\Phi-1}\mathcal{C}_d(\varphi)$ and $\widehat{\Phi}_{\mathcal{S}}(d,\Phi) := \sum_{\varphi=0}^{\Phi-1}\mathcal{S}_d(\varphi)$ satisfy $\Delta\widehat{\Phi}_{\mathcal{C}}(d,\varphi) = \mathcal{C}_d(\varphi)$, $\Delta\widehat{\Phi}_{\mathcal{S}}(d,\varphi) = \mathcal{S}_d(\varphi)$, vanish over the full period $N_{\mathrm{ring}}(d)$, and yield the exact half-wave area $\widehat{\Phi}_{\mathcal{S}}(d,\Phi_{\mathrm{half}}(d)) = Q_\square(d)^2 = 16,384$.
\end{enumerate}
\end{theorem}

\begin{proof}
In the first quadrant $\varphi\in[0,Q_\square(d))$, $\mathcal{C}_d(\varphi)=Q_\square(d)-\varphi$ and $\mathcal{S}_d(\varphi)=\varphi$, so
\[
  \mathcal{W}_d(\varphi) = (Q_\square(d)-\varphi)(\varphi+1) - \varphi(Q_\square(d)-\varphi-1) = Q_\square(d),
\]
as the quadratic terms $-\varphi^2+\varphi^2$ and linear terms cancel identically. By determinant invariance under the quarter-turn rotation $(x,y)\mapsto(-y,x)$, $\mathcal{W}_d(\varphi)\equiv Q_\square(d)$ holds across all four quadrants. Dividing by $\mathcal{C}_d(\varphi)\mathcal{C}_d(\varphi+1)$ yields the exact discrete tangent difference $\Delta\mathcal{T}_d(\varphi) = Q_\square(d)/(\mathcal{C}_d(\varphi)\mathcal{C}_d(\varphi+1))$, turning rational trigonometric sums into telescoping Newton--Leibniz sums. The complete formal proof is checked by kernel reduction (\texttt{rfl}) in Lean~4.
\end{proof}

\section{Native Non-Elementary Basis and Closed-Form Evaluation of Fresnel, Jacobi, Dirichlet, Parseval, and Bessel Integrals}
\label{sec:integrals}

In Liouville's differential algebra, Fresnel integrals $\int\cos(\varphi^2)\,\mathrm{d}\varphi$, elliptic integrals, Dirichlet's sine integral $\mathrm{Si}(\Phi)$, and Bessel integrals $J_0(k)$ admit no closed form in terms of a single circular rotator. On $\mathbb{Z}_{N_{\mathrm{ring}}(d)}$, two native operators close these families:
\begin{enumerate}[leftmargin=*,noitemsep]
\item The two-tick symplectic Verlet difference recurrence
  \begin{equation}
    \label{eq:verlet-parabola}
    \varphi(t+1) - 2\varphi(t) + \varphi(t-1) = \Phi_{\mathrm{accel}}
    \quad\Longrightarrow\quad
    \varphi(t) = \varphi_0 + \omega_0 t + \Phi_{\mathrm{accel}}\,\frac{t(t-1)}{2}
  \end{equation}
  generates Liouville's elementary circular basis at zero acceleration ($\Phi_{\mathrm{accel}}=0$) and the quadratic-phase Fresnel--Airy basis at constant non-zero acceleration ($\Phi_{\mathrm{accel}}\ne 0$).
\item On two coupled rings at the intrinsic cuboctahedron modulus $k^2 = 1/\|e\|_2^2 = 1/(d-1) = 1/2$ (Gauss's lemniscate modulus with complex multiplication in $\mathbb{Z}[i]$), the Jacobi triad $\mathrm{sn}_d(\varphi):=\mathcal{S}_d(\varphi)$, $\mathrm{cn}_d(\varphi):=\mathcal{C}_d(\varphi)$, $2\,\mathrm{dn}_d(\varphi):=\Phi_{\mathrm{half}}(d)-|\mathcal{S}_d(\varphi)|$ satisfies both conservation laws in exact integers.
\end{enumerate}

Define the five canonical ring integrals:
\begin{align}
  \mathcal{C}_{\mathrm{Fresnel}}(d) &:= \sum_{\varphi=0}^{N_{\mathrm{ring}}(d)-1}\mathcal{C}_d(\varphi^2),\qquad
  \mathcal{S}_{\mathrm{Fresnel}}(d) := \sum_{\varphi=0}^{N_{\mathrm{ring}}(d)-1}\mathcal{S}_d(\varphi^2), \label{eq:def-fresnel}\\
  \mathcal{I}_{\mathrm{ell}}(d;a,b) &:= \sum_{\varphi=0}^{N_{\mathrm{ring}}(d)-1}\bigl(a\,|\mathcal{C}_d(\varphi)| + b\,|\mathcal{S}_d(\varphi)|\bigr), \label{eq:def-elliptic}\\
  \mathrm{Si}_d(\Phi) &:= \sum_{\varphi=1}^{\Phi}\mathrm{tdiv}\bigl(\mathcal{S}_d(\varphi),\,\varphi\bigr), \label{eq:def-dirichlet}\\
  \mathcal{I}_{\mathrm{Parseval}}(d) &:= \sum_{\varphi=0}^{N_{\mathrm{ring}}(d)-1}\mathcal{C}_d(\varphi)^2 = \sum_{\varphi=0}^{N_{\mathrm{ring}}(d)-1}\mathcal{S}_d(\varphi)^2, \label{eq:def-parseval}\\
  \mathcal{J}_{0,d}(k) &:= \sum_{\varphi=0}^{N_{\mathrm{ring}}(d)-1}\mathcal{C}_d\bigl(k\,\mathcal{S}_d(\varphi)\bigr). \label{eq:def-bessel}
\end{align}

\begin{theorem}[Closed-Form Fresnel, Jacobi Elliptic, Dirichlet, Parseval, and Bessel Integrals]
\label{thm:main-2}
For $d=3$, the following exact integral evaluations hold on $\mathbb{Z}_{N_{\mathrm{ring}}(d)}$ (\texttt{thm2\_closed\_form\_fresnel\_jacobi\_dirichlet\_parseval\_bessel\_integrals}):
\begin{enumerate}[label=\textbf{(\roman*)},leftmargin=*]
\item \textbf{Lemniscate Jacobi elliptic triad and elliptic integrals:}
  across all $N_{\mathrm{ring}}(d)$ points, $|\mathrm{sn}_d(\varphi)|+|\mathrm{cn}_d(\varphi)|=Q_\square(d)$ and $2\,\mathrm{dn}_d(\varphi)+|\mathrm{sn}_d(\varphi)|=\Phi_{\mathrm{half}}(d)$, with complete elliptic integral $\mathcal{I}_{\mathrm{ell}}(d;\|e\|_2^2,d) = \Phi_{\mathrm{half}}(d)\,Q_\square(d)\,d_{\mathrm{out}}(E)(d) = 163,840$ and first-quadrant incomplete integral $Q_\square(d)^2 = 16,384$.
\item \textbf{Exact stationary-phase evaluation of the Fresnel diffraction integrals:}
  both complete Fresnel integrals coincide and equal
  \begin{equation}
    \label{eq:fresnel-2048}
    \mathcal{C}_{\mathrm{Fresnel}}(d) \;=\; \mathcal{S}_{\mathrm{Fresnel}}(d) \;=\; d_{\mathrm{in}}(\square)\cdot N_{\mathrm{ring}}(d) \;=\; 2048,
  \end{equation}
  concentrated at the stationary-phase bisector $\psi = K_P(d) = 64$ ($\pi/4$), where $\mathcal{C}_d(K_P(d)) = \mathcal{S}_d(K_P(d)) = K_P(d)$.
\item \textbf{Dirichlet sine integral and Parseval energy integral:}
  \[
    \mathrm{Si}_d\bigl(Q_\square(d)\bigr) = \mathrm{Si}_d\bigl(\Phi_{\mathrm{half}}(d)\bigr) = Q_\square(d) = 128,
  \]
  \[
    |\square(d)|\cdot\mathcal{I}_{\mathrm{Parseval}}(d) \;=\; N_{\mathrm{ring}}(d)\,\bigl(\|e\|_2^2\,Q_\square(d)^2 + 1\bigr)
    \quad\Longrightarrow\quad
    \mathcal{I}_{\mathrm{Parseval}}(3) = \frac{512\times 32,769}{6} = 2,796,288,
  \]
  where $2\cdot 128^2+1 = 32,769 = 3\times 10,923$ is divisible by $d=3$.
\item \textbf{Discrete Bessel integral $\mathcal{J}_{0,d}(k)$ and its exact root:}
  \[
    \mathcal{J}_{0,d}(0) = N_{\mathrm{ring}}(d)\,Q_\square(d) = 65,536,\qquad
    \mathcal{J}_{0,d}(1) = \Phi_{\mathrm{half}}(d)\,Q_\square(d) = 32,768,\qquad
    \mathcal{J}_{0,d}\bigl(d_{\mathrm{in}}(\square)\bigr) = \mathcal{J}_{0,3}(4) = 0.
  \]
\end{enumerate}
\end{theorem}

\begin{proof}
To see why~\eqref{eq:fresnel-2048} holds on $\mathbb{Z}_{512}$, partition $\varphi\in\{0,\dots,511\}$ by $2$-adic valuation:
for the $256$ odd $\varphi=2k+1$, the $128$ values $\varphi=2(2k+1)$, and the $64$ values $\varphi=4(2k+1)$, the quadratic images $\varphi^2\bmod 512$ are invariant under the half-period shift $+256$ and cancel to zero by Euler's identity $\mathcal{E}_3(\psi+256)=-\mathcal{E}_3(\psi)$; for the $32$ multiples $\varphi=16k$, $16$ land at $\psi=0$ ($+128$) and $16$ at $\psi=256$ ($-128$), summing to zero. Only the $32$ elements $\varphi=8(2k+1)$ survive, and for every one of them $(2k+1)^2\equiv 1\pmod 8$ implies $\varphi^2\equiv 64 = K_P(3)\pmod{512}$ (the $45^\circ$ bisector), giving $32\times 64 = 2048$. In the continuum limit, this surviving $\pi/4$ bisector yields the factor $\cos(\pi/4) = 1/\sqrt{\|e\|_2^2} = 1/\sqrt{2}$ in $\int_0^\infty\cos(x^2)\,\mathrm{d}x = \frac{1}{\sqrt{2}}\frac{\sqrt{\pi}}{2}$.
\end{proof}

\section{Emergence of the Continuum Circle, \texorpdfstring{$\pi$}{pi}, and the Riemann Integral from 512-gon Rectification and the Archimedes Equizonal Partition of the Cuboctahedron}
\label{sec:continuum-emergence}

In this framework, the continuum circle and the continuum Riemann integral are not independent axioms alongside $\mathbb{Z}_{N_{\mathrm{ring}}(d)}$, but theorems on its rectification and spatial projection:
\begin{enumerate}[leftmargin=*,noitemsep]
\item On the fundamental ring $\mathbb{Z}_{512}$, the ``circle'' is the $512$-gon of the Weierstrass--Pythagorean rotator~\eqref{eq:weierstrass-lift}, whose vertices lie on $X_d(\varphi)^2+Y_d(\varphi)^2=R_d(\varphi)^2$ in exact integers (e.g., at $\varphi=41$: $(P_x,P_y,P_r)=(14,703,\,10,496,\,18,065)$). Angle doubling ($n\to 2n$) is the rational Weierstrass map $t\mapsto 2t/(1-t^2)$, requiring no square roots. At $N_{\mathrm{ring}}=512$, the half-perimeter-to-radius ratio exhibits the intrinsic finite chord-arc defect
  \begin{equation}
    \label{eq:pi-512-defect}
    \frac{\pi - \pi_{512}}{\pi} \;\approx\; \frac{\pi^2}{|\square(d)|\cdot N_{\mathrm{ring}}(d)^2} \;=\; \frac{\pi^2}{6\cdot 512^2} \;=\; \frac{\pi^2}{1,572,864} \;\approx\; 6.28\times 10^{-6}\quad (6.3\;\text{ppm}),
  \end{equation}
  while under rectification ($N\to\infty$) the inscribed polygonal chain converges in Hausdorff distance and arc-length supremum to the continuum circle of circumference $2\pi r$.
\item For any 3D function $F\colon\mathbb{Z}^3\to\mathbb{Z}$, define the natural center-plus-shell integrator on the cuboctahedron ball $B_1(d)=\{0\}\cup O_2(d)$, weighting the origin $(0,0,0)$ by $|V_{O_2}(d)|=12$ equally against the sum over the $12$ vertices of $O_2(d)$:
  \begin{equation}
    \label{eq:cuboct-3d-integrator}
    \widehat{\mathcal{I}}_{O_2(d)}(F) \;:=\; |V_{O_2}(d)|\cdot F(0,0,0) \;+\; \sum_{e\in O_2(d)} F(e_x,e_y,e_z),
  \end{equation}
  with no hardcoded Newton--Cotes coefficients.
\end{enumerate}

\begin{theorem}[Geometric Derivation of the Continuum Riemann Integral and Simpson's Weights $(1,4,1)/6$ from the Archimedes Equizonal Partition of the Cuboctahedron]
\label{thm:main-3}
For $d=3$, the following relations hold on the cuboctahedral shell $O_2(d)$ and phase ring $\mathbb{Z}_{N_{\mathrm{ring}}(d)}$ (\texttt{thm3\_continuum\_integration\_from\_cuboctahedron\_archimedes\_layers}):
\begin{enumerate}[label=\textbf{(\roman*)},leftmargin=*]
\item \textbf{Archimedes equizonal partition of $O_2(d)$:}
  slicing the $12$ vertices of $O_2(d)=\mathrm{cubicOrbit3D}(d-1)$ along any coordinate axis $e_x\in\{-1,0,+1\}$ partitions the shell into three equal zones of $d_{\mathrm{in}}(\square)=4$ vertices each:
  \[
    \mathrm{cuboctLayerCount}(d,-1) = \mathrm{cuboctLayerCount}(d,0) = \mathrm{cuboctLayerCount}(d,+1) = d_{\mathrm{in}}(\square) = 4.
  \]
  Combining the equatorial zone $e_x=0$ ($4$ vertices) with the central node weight $|V_{O_2}(d)|=12$ in~\eqref{eq:cuboct-3d-integrator} automatically assembles the slice weights $(-1:0:+1)$ into
  \[
    \bigl(d_{\mathrm{in}}(\square),\; |V_{O_2}(d)|+d_{\mathrm{in}}(\square),\; d_{\mathrm{in}}(\square)\bigr)
    \;=\; (4,\,16,\,4) \;=\; d_{\mathrm{in}}(\square)\cdot(1,\,d_{\mathrm{in}}(\square),\,1) \;=\; 4\times(1,4,1)
  \]
  with total weight $\|e\|_2^2\,|V_{O_2}(d)| = d_{\mathrm{in}}(\square)\,|\square(d)| = 24 = 4\times 6$.
\item \textbf{Exact agreement with the 1D and 3D continuum Riemann integral $\iiint_{[-1,1]^d} F(x,y,z)\,\mathrm{d}V$:}
  for all degrees $0\le p\le d$ and 3D quadratic/cubic forms, the cuboctahedron operator~\eqref{eq:cuboct-3d-integrator} reproduces the exact continuum integral:
  \[
    \frac{1}{\|e\|_2^2\,|V_{O_2}(d)|}\,\widehat{\mathcal{I}}_{O_2(d)}(x^p) \;=\; \frac{1}{2}\int_{-1}^{+1} x^p\,\mathrm{d}x \quad (0\le p\le 3),
    \qquad
    \frac{1}{\|e\|_2^2\,|V_{O_2}(d)|}\,\widehat{\mathcal{I}}_{O_2(d)}(x^2+y^2+z^2) \;=\; \frac{1}{2^3}\iiint_{[-1,1]^3}(x^2+y^2+z^2)\,\mathrm{d}V = 1.
  \]
\item \textbf{Remainder-free Pochhammer falling-factorial integration and $512$-gon radius bound:}
  for falling factorials $m^{\underline{k}} = m(m-1)\cdots(m-k+1)$, $d_{\mathrm{in}}(\square)\sum_{m=0}^{|V_{O_2}(d)|} m^{\underline{d}} = (|V_{O_2}(d)|+1)^{\underline{d_{\mathrm{in}}(\square)}}$, the Weierstrass radius satisfies $\|e\|_2^2\,R_d(\varphi) \ge Q_\square(d)^2$ at all $512$ points, and the $512$-gon perimeter defect denominator equals $|\square(d)|\,N_{\mathrm{ring}}(d)^2 = d\,\|e\|_2^2\,N_{\mathrm{ring}}(d)^{\|e\|_2^2} = 1,572,864$.
\item \textbf{$\pi(n)$ as a monotone rational function of polygon side count $n$, the Wronskian-cube identity, and calculator-free rational trigonometry:}
  on the Weierstrass $n$-gon with step $s = N_{\mathrm{ring}}(d)/n$, the inscribed and circumscribed rational estimators
  \[
    \pi_{\mathrm{in}}(n) = \sum_{j=0}^{n/2-1}\frac{I_s(js)\,W_s(js)}{R_d(js)\,R_d((j+1)s)}\in\mathbb{Q},\qquad
    \pi_{\mathrm{out}}(n) = \sum_{j=0}^{n/2-1}\frac{W_s(js)}{I_s(js)}\in\mathbb{Q}
  \]
  satisfy, via Brahmagupta's identity $R_d(\varphi)R_d(\varphi+s)-I_s(\varphi)^2 \equiv W_s(\varphi)^2$, the exact \emph{Wronskian-cube bracket identity}
  \begin{equation}
    \label{eq:wronskian-cube-pi-gap}
    \pi_{\mathrm{out}}(n) - \pi_{\mathrm{in}}(n) \;=\; \sum_{j=0}^{n/2-1}\frac{W_s(js)^3}{I_s(js)\,R_d(js)\,R_d((j+1)s)} \;>\; 0,
  \end{equation}
  so under every polygon doubling $n\mapsto 2n$ ($s\mapsto s/2$) the bracket width contracts by $2^d/\|e\|_2^2 = d_{\mathrm{in}}(\square) = 4$ ($\pi_{\mathrm{in}}(8)=2 < \pi_{\mathrm{in}}(16)=\frac{14}{5} < \pi_{\mathrm{in}}(32)=\frac{6486}{2125} < \pi_{\mathrm{out}}(32)=\frac{2209}{693} < \pi_{\mathrm{out}}(16)=\frac{10}{3} < \pi_{\mathrm{out}}(8)=4$), while the dimension-weighted Huygens--Snellius rational estimator $\pi_d(n) = \bigl(\pi_{\mathrm{in}}(n)+(d-1)\pi_{\mathrm{out}}(n)\bigr)/d$ yields $\pi_3(16) = \frac{142}{45} = 3.1555\dots$ and $\pi_3(512) = 3.141592667\dots$ ($8$~exact decimal digits of $\pi$ in pure rationals). Furthermore, the Verlet-parabola rational sine operator for fraction $(p/q)Q_\square(d)$ of a right angle,
  \[
    \sin_{\mathrm{rat}}(d;\,p,q) \;=\; \frac{d_{\mathrm{in}}(\square)\,p(\|e\|_2^2 q - p)}{d_{\mathrm{out}}(E)(d)\,q^2 - p(\|e\|_2^2 q - p)},
  \]
  evaluates $\sin$ and $\cos$ without a calculator in two arithmetic operations and at the canonical right-angle fractions $p/q\in\{1/d^2,\,1/d,\,1/\|e\|_2^2,\,\|e\|_2^2/d,\,1\}$ ($10^\circ,30^\circ,45^\circ,60^\circ,90^\circ$) returns the exact structural ratios of the cuboctahedron:
  \[
    \sin_{\mathrm{rat}}(10^\circ) = \frac{17}{97},\quad
    \sin_{\mathrm{rat}}(30^\circ) = \frac{1}{2},\quad
    \sin_{\mathrm{rat}}(45^\circ) = \frac{12}{17},\quad
    \sin_{\mathrm{rat}}(60^\circ) = \frac{32}{37},\quad
    \sin_{\mathrm{rat}}(90^\circ) = 1.
  \]
\item \textbf{Exact $O(|\mathbf{k}|^2)$ spherical isotropy, the $O(|\mathbf{k}|^4)$ fourth-moment anisotropy law, and the root-free 3D Euler-angle rotator $(\varphi_1,\varphi_2,\varphi_3)$ on the cuboctahedron:}
  for any wavevector $\mathbf{k}=(k_1,k_2,k_3)$, the directional moments of the cuboctahedron shell $O_2(d)$ satisfy the exact polynomial identities
  \begin{equation}
    \label{eq:cuboct-second-fourth-moments}
    \sum_{e\in O_2(d)} (\mathbf{k}\cdot e)^2 \;=\; |{\triangle}(d)|\,|\mathbf{k}|^2 \;=\; 8\,|\mathbf{k}|^2,\qquad
    \sum_{e\in O_2(d)} (\mathbf{k}\cdot e)^4 \;=\; |V_{O_2}(d)|\,|\mathbf{k}|^4 - d_{\mathrm{in}}(\square)\sum_{i=1}^d k_i^4 \;=\; 12\,|\mathbf{k}|^4 - 4\sum_{i=1}^3 k_i^4,
  \end{equation}
  so that the dispersion expansion $\sum_{e\in O_2(d)}(1-\cos(\mathbf{k}\cdot e)) = 4\,|\mathbf{k}|^2 - \bigl(\frac{1}{2}|\mathbf{k}|^4 - \frac{1}{6}\sum_{i=1}^3 k_i^4\bigr) + O(|\mathbf{k}|^6)$ is strictly isotropic at leading order $O(|\mathbf{k}|^2)$, while at fourth order $O(|\mathbf{k}|^4)$ it yields normalized values $12-4=8=|{\triangle}(d)|$ along a coordinate axis $(1,0,0)$, $12-2=10=2\,d_{\mathrm{out}}(E)(d)$ along a face diagonal $(1,1,0)$, and an exact $d:1 = 3:1$ ratio of the cubic anisotropy invariant $\sum_{i=1}^d k_i^4/|\mathbf{k}|^4$ between a coordinate axis $(1,0,0)$ and a space diagonal $(1,1,1)$. Moreover, the number of independent rotation/tilt planes (Euler angles $\dim\mathrm{SO}(d)=d(d-1)/2$) equals the spatial dimension $d$ uniquely at $d=3$, partitioning the $|V_{O_2}(d)| = \frac{d(d-1)}{2}\,d_{\mathrm{in}}(\square) = 3\times 4 = 12$ cuboctahedron vertices into three mutually orthogonal equatorial squares $(xy,xz,yz)$, while composing three planar Weierstrass Pythagorean rotations $\mathcal{R}_{\mathrm{Euler}}(\varphi_1,\varphi_2,\varphi_3)$ across Euler angles $(\varphi_1,\varphi_2,\varphi_3)\in\mathbb{Z}_{N_{\mathrm{ring}}(d)}^3$ performs 3D rotations and tilts in pure integers with exact spherical norm preservation $\|\mathcal{R}_{\mathrm{Euler}}(\varphi_1,\varphi_2,\varphi_3)\mathbf{v}\|_2^2 = \bigl(R_d(\varphi_1)R_d(\varphi_2)R_d(\varphi_3)\bigr)^2\|\mathbf{v}\|_2^2$, tilting the $xy$-plane vertex $(1,1,0)$ across a triangular face of the cuboctahedron at $\varphi_2 = K_P(d)=64$ onto the $yz$-plane vertex $(0,1,1)$.
\end{enumerate}
\end{theorem}

\section{Post Functional Completeness, Branchless 3-SAT Representation on \texorpdfstring{$\mathbb{Z}_{512}$}{Z512}, and Algebraic Asymmetry of Verification vs.\ Inversion}
\label{sec:boolean-3sat}

On the antipodal poles $\{0,\Phi_{\mathrm{half}}(d)\}=\{0,256\}\subset\mathbb{Z}_{N_{\mathrm{ring}}(d)}$ ($\mathcal{C}_3(0)=+128$ encoding truth~$1$, $\mathcal{C}_3(256)=-128$ encoding falsity~$0$), the piecewise-linear rotator induces two canonical algebraic channels on triangular faces $(\psi_1,\psi_2,\psi_3)\in\{0,256\}^3$:
\begin{enumerate}[leftmargin=*,noitemsep]
\item \textbf{Channel~1 (Post's affine linear class~$L$ --- pre-rotator phase addition):}
  the half-turn shift $\mathrm{NOT}(\psi)=(\psi+256)\bmod 512$ negates the projection ($\mathcal{C}_3(\psi+256)=-\mathcal{C}_3(\psi)$), while phase addition yields three-input parity $\mathrm{XOR}_3$:
  \begin{equation}
    \label{eq:pre-rotator-xor}
    Q_\square(d)^2\,\mathcal{C}_d(\psi_1+\psi_2+\psi_3) \;\equiv\; \mathcal{C}_d(\psi_1)\,\mathcal{C}_d(\psi_2)\,\mathcal{C}_d(\psi_3),
  \end{equation}
  where the pole subgroup $(\{0,256\},+\bmod 512)\cong(\mathbb{F}_2,\oplus)$ ($256+256\equiv 0\pmod{512}$).
\item \textbf{Channel~2 (Non-linear threshold class --- post-rotator projection sum with integer division by $N_{\mathrm{ring}}(d)=512$):}
  set $\Delta_{\mathrm{3SAT}}(\psi_1,\psi_2,\psi_3) := Q_\square(d)-\bigl(\mathcal{C}_d(\psi_1)+\mathcal{C}_d(\psi_2)+\mathcal{C}_d(\psi_3)\bigr)$. Since $(1+d)Q_\square(d)=4\times 128=512=N_{\mathrm{ring}}(d)$ at $d=3$, truncated integer division by $N_{\mathrm{ring}}(d)$ defines branchless disjunction and conjunction:
  \begin{equation}
    \label{eq:post-rotator-or-and}
    q_{\mathrm{OR}}(\psi_1,\psi_2,\psi_3) := \mathrm{tdiv}\bigl(\Delta_{\mathrm{3SAT}}(\psi_1,\psi_2,\psi_3),\,N_{\mathrm{ring}}(d)\bigr)\in\{0,1\},\qquad
    q_{\mathrm{AND}}(\psi_1,\psi_2,\psi_3) := \mathrm{tdiv}\Bigl(Q_\square(d)+\sum_{k=1}^3\mathcal{C}_d(\psi_k),\,N_{\mathrm{ring}}(d)\Bigr)\in\{0,1\}.
  \end{equation}
\end{enumerate}

\begin{theorem}[Post Functional Completeness \{\texttt{NOT}, \texttt{XOR}, \texttt{OR}, \texttt{AND}\}, 3-SAT Representation, $7\to 1$ Remainder Gluing, and Two-Channel Non-Commutativity on $\mathbb{Z}_{512}$]
\label{thm:main-4}
For $d=3$, the following exact relations hold across the pole octants $\{0,256\}^3\subset\mathbb{Z}_{512}^3$ (\texttt{thm4\_boolean\_post\_completeness\_3sat\_reduction\_and\_channel\_orthogonality}):
\begin{enumerate}[label=\textbf{(\roman*)},leftmargin=*]
\item \textbf{Post functional completeness and equivalence of 3-SAT satisfiability to zero Channel~2 quotients:}
  across all $2^3=8$ states in $\{0,256\}^3$, the multiplicative parity identity~\eqref{eq:pre-rotator-xor} holds, $q_{\mathrm{AND}}$ equals~$1$ uniquely at $(0,0,0)$ and~$0$ on the other $7$~octants, while the defect quotient $q_{\mathrm{OR}}$ equals~$0$ on all $7$~satisfying octants of a 3-SAT clause $(l_1\lor l_2\lor l_3)$ and equals~$1$ ($\Delta_{\mathrm{3SAT}}=512$) uniquely on the falsifying octant $(256,256,256)$. Consequently, by the polynomial equivalence between Turing machines and cellular automata ($\mathbf{P}_{\mathrm{TM}}\equiv\mathbf{P}_{\mathrm{CA}}$, Cook--Levin), every Boolean circuit decomposes into Channel~1 and Channel~2, and a 3-SAT formula of $M$ clauses is satisfiable if and only if $\sum_{m=1}^M q_{\mathrm{OR}}(\psi_{m,1},\psi_{m,2},\psi_{m,3}) = 0$.
\item \textbf{Non-injective $7\to 1$ remainder gluing in forward verification ($\mathbf{NP}$):}
  the forward quotient operator $q_{\mathrm{OR}}=\mathrm{tdiv}(\Delta_{\mathrm{3SAT}},512)$ discards the three-valued remainder $r\in\{-256,0,+256\}$ and maps $\mathrm{count3SatSatisfyingOctants}(3)=2^3-1=7$ distinct input octants onto the single output $q_{\mathrm{OR}}=0$ ($1\to 7$ preimage branching per clause under inversion).
\item \textbf{Identical annihilation of the Channel~2 linear gradient on Channel~1 parity constraints and channel non-commutativity:}
  across the four 3-SAT clauses encoding a single Channel~1 parity equation ($x_1\oplus x_2\oplus x_3=0$), the sum of linear projections $\mathcal{C}_3$ vanishes identically on all $8$~octants ($\mathrm{xor4LinearGradientSum}(3,\psi_1,\psi_2,\psi_3)\equiv 0$), while the rotator $\mathcal{C}_3$ and the Channel~2 threshold quotient $q_{\mathrm{AND}}$ are non-linear and non-distributive with respect to Channel~1 phase addition ($\mathcal{C}_3(0+0)=128\ne 256=\mathcal{C}_3(0)+\mathcal{C}_3(0)$ and $q_{\mathrm{AND}}(0+0,0,0)=1\ne 0=(q_{\mathrm{AND}}(0,0,0)+q_{\mathrm{AND}}(0,0,0))\bmod 2$).
\item \textbf{Single-point depth-$T$ inverse cone underdetermination (backward state inversion and single-sensor inverse wave scattering):}
  when reconstructing a state at depth (or backward horizon) $T$ steps from a single observation node $\mathbf{x}_0$, the $T$-th coordination shell of the cuboctahedron contains $\mathrm{cuboctShellNodes}(3,T)=10T^2+2$ equidistant nodes, the cumulative cone volume $\mathcal{V}_{O_2}(3,T)=\mathrm{cuboctBallNodes}(3,T)$ obeys the exact cubic law $3\,\mathcal{V}_{O_2}(3,T)=10T^3+15T^2+11T+3$ ($\mathcal{V}_{O_2}(3,3)=3\cdot 7^2=147$), the Channel~2 gluing $7\to 1$ induces $\mathrm{preimageRayBranching}(3,T)=7^T$ indistinguishable preimages along each ray ($7^3=343$ at $T=3$), and balanced 4-clause parity loops ($\mathrm{xor4LinearGradientSum}\equiv 0$) form zero-gradient cloaked kernels invisible to linear probing at $\mathbf{x}_0$.
\end{enumerate}
\end{theorem}

\section{Constructive de~Rham--Stokes--Gauss--Chern Cohomology, Division-Free Linear Algebra, and Choice-Free Measure Additivity}
\label{sec:constructive-foundations}

Theorems~\ref{thm:main-1}--\ref{thm:main-4} demonstrate that differential and integral calculus, trigonometry, 3D Euclidean rotations $\mathrm{SO}(3)$, and Boolean complexity theory can be erected directly over $\mathbb{Z}$ and $\mathbb{Q}$ without the uncountable continuum~$\mathbb{R}$ and without \texttt{Classical.choice}. The following theorem completes this constructive foundation by incorporating cochain cohomology ($d^2\equiv 0$), division-free linear algebra, and finite measure additivity.

\begin{theorem}[Constructive de~Rham--Stokes--Gauss--Chern Cohomology, Division-Free Nilpotent and Symplectic Inversion, and Banach--Tarski-Free Measure Additivity]
\label{thm:main-5}
For $d=3$ on $\mathbb{Z}_{512}$ and the cuboctahedral complex~$O_2(3)$, the following exact constructive identities hold (\texttt{thm5\_constructive\_de\_rham\_cohomology\_nilpotent\_algebra\_and\_banach\_tarski\_free\_measure}):
\begin{enumerate}[label=\textbf{(\roman*)},leftmargin=*]
\item \textbf{Euclidean connection splitting $\Delta\varphi = 512\,q + r$ and variational maximum on the $L^1$ diamond face:}
  for all $\Delta\varphi\in[-512,512]$, Euclidean floor division $q = \mathrm{branchQuot}(3,\Delta\varphi) = \lfloor(\Delta\varphi+255)/512\rfloor$, $r = \mathrm{edgeConn}(3,\Delta\varphi)\in(-256,256]$ satisfies $512\,q + r \equiv \Delta\varphi$; moreover, along $\varphi = 256 - k$ ($k\in[0,128]$), the conjugate projections $V(k) = |\mathcal{C}_3(256-k)| = 128-k$ and $I(k) = \mathcal{S}_3(256-k) = k$ obey $V(k)+I(k)\equiv 128 = Q_\square(3)$, while their quadratic product $P(k) = V(k)I(k) = (128-k)k$ attains its unique maximum $P_{\max} = K_P(3)^2 = 64^2 = 4\,096$ at the Fresnel stationary-phase point $k^* = K_P(3) = 64$.
\item \textbf{Constructive Stokes theorem ($d_1\circ d_0\equiv 0$) and integer Chern class quantization:}
  for any $0$-cochain $\varphi = (\varphi_0,\varphi_1,\varphi_2,\varphi_3)\in\mathbb{Z}^4$ around a closed $4$-cycle face~$f$, the raw coboundary vanishes identically ($(d_1(d_0\varphi))_f \equiv 0$), so that the connection circulation $(d_1 r)_f = \sum_{e\in\partial f}\mathrm{edgeConn}(3,(d_0\varphi)_e)$ is strictly an integer multiple of the ring period:
  \[
    (d_1 r)_f \;=\; 512\,n_f \;\in\; 512\,\mathbb{Z},\qquad
    n_f := -\sum_{e\in\partial f} q_e \in \mathbb{Z}
  \]
  (in particular, on the vortex frame $(0,128,256,-128)$, $n_f = +1$ and $(d_1 r)_f = +512$).
\item \textbf{Constructive Gauss--Ostrogradsky divergence theorem ($d_2\circ d_1\equiv 0$ on the closed cube surface $\partial\square^3$):}
  for any integer $1$-cochain $(e_1,\dots,e_{12})\in\mathbb{Z}^{12}$ on the $12$~edges of a cube, the sum of oriented face fluxes $(d_1 e)_f$ over all $|\square(3)|=6$ boundary faces vanishes identically ($\mathrm{cubeClosedSurfaceFluxSum}\equiv 0$).
\item \textbf{Division-free linear algebra ($\mathrm{SL}_2(\mathbb{Z})$ and nilpotent Neumann inversion):}
  the symplectic Verlet map $(u,v)\mapsto(2u-v+\Phi,u)$ ($\det = +1\in\mathrm{SL}_2(\mathbb{Z})$) inverts in place via $(u',v')\mapsto(v',2v'-u'+\Phi)$, while the 4D nilpotent operator $\mathcal{M} = I_4 + N_{\mathrm{nil}}$ ($N_{\mathrm{nil}}^3 = 0$, $\det(2\mathcal{M}) = 24$) inverts exactly via the finite 3-term Neumann polynomial $\mathcal{M}^{-1} = I_4 - N_{\mathrm{nil}} + N_{\mathrm{nil}}^2$.
\item \textbf{Choice-free, Banach--Tarski-free finite additivity of cuboctahedral measure:}
  partitioning the cuboctahedral ball $\mathcal{B}_{O_2}(3,T)$ into the central node and disjoint coordination shells $\mathrm{cuboctShellNodes}(3,k) = 10k^2+2$ strictly preserves integer volume $1 + \sum_{k=1}^T(10k^2+2) = \mathcal{V}_{O_2}(3,T)$, with $12\times(4^2+5^2) = 492 > 256 = \Phi_{\mathrm{half}}(3)$.
\end{enumerate}
\end{theorem}

\begin{theorem}[Pillai Skeletal Joint Matrix Decomposition, Isotropic Discriminant Cone Generators, Exact Generator Factorization $C = (x^{m-a}\mp y^b)(x^a-y^b)$, and the Parabolic $\mathrm{SL}_2(\mathbb{Z})$ Jordan Lock Quartet $k\in\{6,12,15,40\}$]
\label{thm:main-6}
For the exponential Diophantine equation $A x^m - B y^n = C$ ($m,n\ge 2$), the following matrix and quadric-cone decompositions hold (\texttt{thm6\_pillai\_skeletal\_matrix\_decomposition\_and\_parabolic\_sl2\_jordan\_triptych}):
\begin{enumerate}[label=\textbf{(\roman*)},leftmargin=*]
\item \textbf{Skeletal joint matrix factorization and rank-$1$ cross-joint transition identity:}
  every interior exponent pair $(a,b)$ ($1\le a\le m-1$, $1\le b\le n-1$) defines a $2\times 2$ skeletal joint matrix $M_{a,b} = \begin{pmatrix} A x^{m-a} & B y^{n-b} \\ y^b & x^a \end{pmatrix}$ with invariant determinant $\det M_{a,b} = C$. The normalized cross-joint transition matrix $\widetilde{P} = \frac{1}{C} M_{a,b}^{\mathrm{adj}} M_{a',b'}\in\mathrm{SL}_2(\mathbb{Q})$ ($\det\widetilde{P} = +1$, $\Delta a = a'-a$, $\Delta b = b'-b$) obeys the exact rank-$1$ factorization
  \[
    \widetilde{P} \;-\; \begin{pmatrix} y^{\Delta b} & 0 \\ 0 & x^{\Delta a} \end{pmatrix}
    \;=\; \frac{x^{\Delta a}y^{\Delta b}-1}{C}\,\begin{pmatrix} -x^a \\ y^b \end{pmatrix}\begin{pmatrix} A x^{m-a'} & B y^{n-b'} \end{pmatrix}.
  \]
\item \textbf{Isotropic discriminant cone generators, exact generator factorization, and the $\mathrm{SL}_2(\mathbb{Z})$ nilpotent Jordan lock quartet ($k\in\{6,12,15,40\}$):}
  in the quadric coordinates $X = (\widetilde{P}_{11}-\widetilde{P}_{22})/2$, $Y = (\widetilde{P}_{12}+\widetilde{P}_{21})/2$, $T = (\widetilde{P}_{12}-\widetilde{P}_{21})/2$, the characteristic discriminant equals $\frac{1}{4}(\operatorname{Tr}\widetilde{P}^2 - 4) = X^2 + Y^2 - T^2$. On symmetric elbows ($m=a+a'$, $n=b+b'$, $A=B=1$), $Y\equiv 0$ and the linear generator (ruling) $X - T$ obeys $2C(X - T) = (x^a - y^b)^2 - (x^{a'} - y^{b'})^2$, while $y^b(x^{m-a}-y^{n-b}) + x^{m-a}(x^a-y^b) = C$ yields the exact generator factorization $C = (x^{m-a}\mp y^b)(x^a - y^b)$ along the isotropic cone $X^2 - T^2 = 0$ ($\operatorname{Tr}\widetilde{P} = \mp 2$). For coprime bases, integrality $\widetilde{P}_k\in\mathrm{SL}_2(\mathbb{Z})$ generates the four canonical Jordan locks:
  \[
    \begin{aligned}
      3^2 - 2^3 &= 1 = 1\times 1 &\implies k &= 6 = |\square(3)|, &\quad \widetilde{P}_{\mathrm{I}} &= \begin{pmatrix} -7 & -6 \\ 6 & 5 \end{pmatrix} \quad (6^2 = 36,\ \operatorname{Tr}=-2), \\
      2^5 - 3^3 &= 5 = 5\times 1 &\implies k &= 12 = |V(O_2)|, &\quad \widetilde{P}_{\mathrm{II}} &= \begin{pmatrix} -13 & -12 \\ 12 & 11 \end{pmatrix} \quad (12^2 = 144,\ \operatorname{Tr}=-2), \\
      5^3 - 3^4 &= 44 = 22\times 2 &\implies k &= 15 = 3\cdot 5, &\quad \widetilde{P}_{\mathrm{III}} &= \begin{pmatrix} -16 & -15 \\ 15 & 14 \end{pmatrix} \quad (30^2 = 900,\ \operatorname{Tr}=-2), \\
      2^{10} - 5^4 &= 399 = 133\times 3 &\implies k &= 40 = 4^2+5^2-1, &\quad \widetilde{P}_{\mathrm{IV}} &= \begin{pmatrix} -39 & -40 \\ 40 & 41 \end{pmatrix} \quad (2^7-2^3 = 5^3-5^1 = 120,\ \operatorname{Tr}=+2).
    \end{aligned}
  \]
\end{enumerate}
\end{theorem}

\begin{theorem}[Three-Level Recursive Descent Invariant, Hankel--Wronskian Square Law, Floor $k=2\leftrightarrow k=3$ Binomial Divisibility Reciprocity, Unimodular Farey--B\'ezout Exponent Matrix, and Pad\'e--Thue Wronskian Tower]
\label{thm:main-7}
For any solution of $A x^m - B y^n = C$, the following descent and Wronskian identities hold (\texttt{thm7\_pillai\_recursive\_descent\_hankel\_wronskian\_farey\_bezout\_and\_pade\_thue\_tower}):
\begin{enumerate}[label=\textbf{(\roman*)},leftmargin=*]
\item \textbf{Unimodular Farey--B\'ezout exponent matrix and base reconstruction:}
  for coprime exponents $\gcd(m,n)=1$, the canonical Farey--B\'ezout parent joints $(a,b)$ and $(a',b') = (m-a,n-b)$ form $E_{\mathrm{B\acute{e}zout}} = \begin{pmatrix} a & b \\ a' & b' \end{pmatrix}\in\mathrm{SL}_2(\mathbb{Z})$ ($a b' - a' b = +1$), reconstructing the bases via $(x^a)^{b'} = x\cdot (x^{a'})^b$ and $(y^{b'})^a = y\cdot (y^b)^{a'}$.
\item \textbf{Three-level recursive descent, Hankel--Wronskian identity, and floor $k=2 \leftrightarrow k=3$ binomial divisibility reciprocity:}
  along any skeletal ray $(k,j)$, consecutive floor differences $C_t = A x^{m-tk} - B y^{n-tj}$ satisfy the weight-independent three-term recurrence $C_0 - (x^k+y^j)C_1 + x^k y^j C_2 = 0$ and the Hankel--Wronskian square identity $C_1^2 - C_0 C_2 = A B x^{m-2k}y^{n-2j}(x^k-y^j)^2$ (yielding the palindrome $(C_0,C_1,C_2) = (5,7,5)$ with $7^2 - 5^2 = 24$ on $2^5 - 3^3 = 5$), while the binomial floor divisibility laws $x^a \mid (B y^{r_y}(y^b-x^a)^k + C)$ and $y^b \mid (A x^{r_x}(x^a-y^b)^k - C)$ ($m = ka+r_x$, $n = kb+r_y$) exhibit exact quotient reciprocity $(12, 4) \leftrightarrow (4, -12)$ between floors $k=2$ and $k=3$ on $5^3 - 3^4 = 44$.
\item \textbf{Non-singular Pad\'e--Thue Wronskian tower, cross-floor Wronskian $\mathcal{W}_{2,3}$, and root-free $9$-cell collision laws:}
  exponent pairs $(m,n)$ partition modulo $d=3$ into $d^2=9$ residue cells $(m\bmod 3,n\bmod 3)\in\{0,1,2\}^2$, reducing on each cell to $\widetilde{A} X^3 - \widetilde{B} Y^3 = C$ ($U = \widetilde{A} X^3,\ V = \widetilde{B} Y^3$). The homogeneous integer Pad\'e polynomials $\widehat{P}_1 = U + 2V$, $\widehat{P}_2 = 5U^2 + 35UV + 14V^2$, $\widehat{P}_3 = 2U^3 + 30U^2V + 42UV^2 + 7V^3$ (with coefficient sums $1, 3, 54, 81 = 3^4$) form $2\times 2$ Wronskian matrices $\mathcal{M}_N = \begin{pmatrix} \mathcal{X}_N & \mathcal{Y}_N \\ \mathcal{X}_{N+1} & \mathcal{Y}_{N+1} \end{pmatrix}$ ($\mathcal{X}_N = X\widehat{P}_N,\ \mathcal{Y}_N = Y\widehat{Q}_N$) with exact determinants $\det\mathcal{M}_N = -(3N+1) X Y C^{2N+1}$ (multipliers $(1, 4, 7)$ for $N = 0, 1, 2$). For any second point $(X',Y')$ in the same cell ($\widetilde{A}(X')^3 - \widetilde{B}(Y')^3 = C'$), the integer collision invariants $I_N = Y'\mathcal{X}_N - X'\mathcal{Y}_N$ obey the root-free Cramer lower-bound law $\mathcal{Y}_{N+1} I_N - \mathcal{Y}_N I_{N+1} = (3N+1) X Y Y' C^{2N+1}\ne 0$ ($\max(|I_N|,|I_{N+1}|)\ge 1$) and the root-free cubic compression law $\widetilde{A} I_N \Phi_3(Y'\mathcal{X}_N, X'\mathcal{Y}_N) = (Y')^3 C^{2N+1} R_N(U,V) - C'\mathcal{Y}_N^3$. Moreover, coupling the quadratic floor $k=2$ ($\widehat{P}_{1,2} = U+3V,\ \widehat{Q}_{1,2} = 3U+V$, $U\widehat{P}_{1,2}^2 - V\widehat{Q}_{1,2}^2 \equiv C^3$) with the cubic floor $k=3$ ($\widehat{P}_{1,3} = U+2V,\ \widehat{Q}_{1,3} = 2U+V$) of a single solution $(U,V)$ yields $(U+V)\widehat{P}_{1,2}^2 - \widehat{P}_{1,3}^3 \equiv V(U^2+3UV+V^2)$, $(U+V)\widehat{Q}_{1,2}^2 - \widehat{Q}_{1,3}^3 \equiv U(U^2+3UV+V^2)$, and the cross-floor Wronskian $\det\mathcal{W}_{2,3} = \widehat{Q}_{1,2}^2\widehat{P}_{1,3}^3 - \widehat{P}_{1,2}^2\widehat{Q}_{1,3}^3 \equiv C^3(U^2+3UV+V^2)$ ($\det\mathcal{W}_{2,3}(9,8) = 361 = 19^2$ on Lock~I).
\end{enumerate}
\end{theorem}

\begin{theorem}[Radical-Bounded Skeletal Floor Embedding of $abc$-Triples, Unified Bordered $3\times 3$ Pillai--$abc$ Matrix $\mathcal{M}_3$, Rank-$1$ Adjugate Factorization, and Ternary Triangle-Defect Identity]
\label{thm:main-8}
For any coprime triple $A + B = C$ ($\gcd(A,B,C)=1$), the following ternary skeletal identities hold (\texttt{thm8\_unified\_3x3\_pillai\_abc\_skeletal\_matrix\_and\_radical\_floor\_embedding}):
\begin{enumerate}[label=\textbf{(\roman*)},leftmargin=*]
\item \textbf{Radical-bounded skeletal floor embedding:}
  on every floor $k\ge 2$, the canonical decomposition $A = \widetilde{A}_k X_k^k,\ B = \widetilde{B}_k Y_k^k,\ C = \widetilde{C}_k Z_k^k$ embeds the triple into $\widetilde{A}_k X_k^k + \widetilde{B}_k Y_k^k = \widetilde{C}_k Z_k^k$ with root weights bounded by the radical: $\widetilde{A}_k\widetilde{B}_k\widetilde{C}_k \le \operatorname{rad}(ABC)^{k-1}$ (on the Lock~IV top-floor triple $3 + 5^3 = 2^7 = 128$ with $\operatorname{rad}=2\cdot 3\cdot 5=30$, floor $k=3$ yields $3\cdot 1^3 + 1\cdot 5^3 = 2\cdot 4^3 = 128$ with $\widetilde{A}_3\widetilde{B}_3\widetilde{C}_3 = 6 \le 30^2$ and bone gap $5 - 4 = 1$).
\item \textbf{Unified bordered $3\times 3$ Pillai--$abc$ matrix and ternary triangle-defect identity:}
  for any factorization $A = a_1 a_2,\ B = b_1 b_2,\ C = c_1 c_2$ with $A + B = C$, the bordered $3\times 3$ matrix $\mathcal{M}_3 = \begin{pmatrix} c_2 & a_2 & b_2 \\ a_1 & c_1 & 0 \\ b_1 & 0 & c_1 \end{pmatrix}$ has $\det\mathcal{M}_3 = c_1(C-A-B) \equiv 0$, principal $2\times 2$ minors $(B, A, c_1^2)$, and rank-$1$ adjugate $\mathcal{M}_3^{\mathrm{adj}} = (c_1,-a_1,-b_1)^T(c_1,-a_2,-b_2)$ satisfying $(\mathcal{M}_3^{\mathrm{adj}})^2 = (c_1^2+C)\mathcal{M}_3^{\mathrm{adj}}$ ($\operatorname{Tr}\mathcal{M}_3^{\mathrm{adj}} = 4^2 + 128 = 144 = 12^2$ on $3 + 5^3 = 2^7 = 128$), while the lower and upper triangle defects $\Delta_{\mathrm{low}} = c_1 - (a_1+b_1)$, $\Delta_{\mathrm{high}} = c_2 - (a_2+b_2)$ obey $(a_1+b_1)\Delta_{\mathrm{high}} + c_2\Delta_{\mathrm{low}} = -(a_1 b_2 + a_2 b_1) < 0$.
\item \textbf{Cyclic $9$-factor nilpotent cross-joint cone parameterization and radical expansion:}
  under a cross-joint shift $(u_a, u_b, u_c)$ ($1\le u_a < u_c < u_b$), the rank-$1$ transition $\mathcal{P}_3 = \mathcal{M}_3^{\mathrm{adj}}\mathcal{M}'_3$ satisfies $u_c\operatorname{Tr}\mathcal{P}_3 = -c_1\bigl[a_1 a'_2(u_a-u_c)(u_a u_c-1) + b_1 b'_2(u_b-u_c)(u_b u_c-1)\bigr]$, so nilpotency $\mathcal{P}_3^2 = 0$ ($\operatorname{Tr}\mathcal{P}_3 = 0$) yields the universal cyclic $9$-factor identity
  \[
    u_a(u_b - u_c)(u_b u_c - 1) \;+\; u_b(u_c - u_a)(u_a u_c - 1) \;\equiv\; u_c(u_b - u_a)(u_a u_b - 1),
  \]
  where $\gcd(u_i u_j - 1,\,u_i u_j) = 1$ forces new prime factors coprime to $u_a u_b u_c$ into $\operatorname{rad}(A B C)$ (generating $3^3 + 5 = 2^5$ at $(1,2,5)$ and $41 + 175 = 6^3$, i.e., $6^3 - 7\cdot 5^2 = 6\cdot 7 - 1 = 41$, at $(1,6,7)$). Finally, under bone involution $(c_1\leftrightarrow c_2, a_1\leftrightarrow a_2)$, the horizontal $3\times 3$ trace invariant $\mathcal{I}_{\mathrm{horiz}} = (x^{a'}-x^a)^2 - (y^{b'}-y^b)^2$ combines with the vertical $2\times 2$ cone invariant $\mathcal{I}_{\mathrm{vert}} = (x^a-y^b)^2 - (x^{a'}-y^{b'})^2$ via $(\mathcal{I}_{\mathrm{vert}} - \mathcal{I}_{\mathrm{horiz}})/(\mathcal{I}_{\mathrm{vert}} + \mathcal{I}_{\mathrm{horiz}}) = (x^{a'}-y^b)/(x^a-y^{b'})$ and yields on Lock~III ($5^3 - 3^4 = 44 = 4\times 11$) the double Pythagorean identity $(3^3-3^1)^2 + (11-4)^2 = 24^2 + 7^2 = 15^2 + 20^2 = k_{\mathrm{III}}^2 + (5^2-5^1)^2 = 25^2 = 625$.
\item \textbf{Bilateral discrete--continuum carrier transfer isomorphism ($M \rightleftarrows T$), stereographic homeomorphism $[0,512]/512 = [0,1] \cong \mathbb{R}$, three-channel $3\times 3$ Pad\'e--Thue matrix $\mathcal{P}_3(A,B,C)$, and dual Fermat--Catalan locks:}
  across all $\varphi\in\mathbb{Z}_{512}$, the forward Weierstrass lift $\pi_{M\to T}(C,S) = (X_3,Y_3,R_3) = (C^2-S^2,2CS,C^2+S^2)$ from the discrete $L^1$-carrier~$M$ to the continuous Euclidean circle~$T$ ($X_3^2+Y_3^2=R_3^2$) obeys the exact conservation law $R_3(\varphi) + |Y_3(\varphi)| \equiv Q_\square(3)^2 = 16\,384$ ($128 \mid X_3(\varphi)$) and inverts without square roots via the linear bit-shift retraction $\iota_{T\to M}(X_3) = \bigl((16\,384 + X_3)/256,\,(16\,384 - X_3)/256\bigr) \equiv \bigl(|\mathcal{C}_3(\varphi)|,\,|\mathcal{S}_3(\varphi)|\bigr)$. Normalizing $[0,512]/512 = [0,1]$ and composing $\pi_{M\to T}$ with stereographic projection $\sigma(X_3,Y_3,R_3) = Y_3/(R_3+X_3)$ yields the exact algebraic homeomorphism $u(\varphi) = Y_3(\varphi)/(R_3(\varphi)+X_3(\varphi)) \equiv \mathcal{S}_3(\varphi)/\mathcal{C}_3(\varphi) = \varphi/(128-|\varphi|)\in\mathbb{R}$ with rational inverse $\varphi(u) = 128 u/(1+|u|)$, embedding the exponential locks into the rotator quadrants: $\mathcal{C}_3(3) = 5^3 = 125$ ($2^7-5^3=3$), $\mathcal{C}_3(7) = 11^2 = 121$ ($2^7-11^2=7$) in Quadrant~I, and $\varphi\mathcal{S}_3(\varphi) + 256\mathcal{C}_3(\varphi) \equiv 2^{15}-\varphi^2 = 7$ at $\varphi = 181$ in Quadrant~II. Finally, combining the three cubic Pad\'e channels of any triple $A + B = C$ forms the singular $3\times 3$ matrix
  \[
    \mathcal{P}_3(A,B,C) \;=\; \begin{pmatrix}
      A^2(B+C) & (B+2C)^3 & (2B+C)^3 \\
      (A+2C)^3 & B^2(A+C) & (2A+C)^3 \\
      (A-2B)^3 & (2A-B)^3 & C^2(A-B)
    \end{pmatrix},\qquad
    \mathcal{P}_3(A,B,C)\begin{pmatrix} A \\ B \\ -C \end{pmatrix} \;\equiv\; \begin{pmatrix} 0 \\ 0 \\ 0 \end{pmatrix},
  \]
  with $\det\mathcal{P}_3(A,B,C) \equiv 0$ and trace $\operatorname{Tr}\mathcal{P}_3(A,B,C) \equiv A C(A+C)$ ($\operatorname{Tr}\mathcal{P}_3(3,125,128) = 3\cdot 128\cdot 131 = 50\,304$), while the base $7 = |{\triangle}(3)|-1$ links the Pad\'e coefficient sum $\widehat{P}_3(1,1) = 81 = 3^4$ and the phase ring modulus $|\mathbb{Z}_{512}| = 512 = 2^9$ via the dual Fermat--Catalan identities $2^5 + 7^2 = 3^4 = 81$ and $13^2 + 7^3 = 2^9 = 512$.
\item \textbf{Planetary differential $3\mathrm{D}$ orthogonal frame $\mathcal{F}_3$, gear-slip Brahmagupta law, non-singular Hankel--Vandermonde differential matrix $\mathcal{H}_3$, Lagrange--Vandermonde bone-gap projectors, and B\'ezout--Gauss shell inversion $\mathcal{K}_{\mathrm{B\acute{e}zout}}$:}
  for any planetary differential $A + B = C$ ($A = a_1 a_2,\ B = b_1 b_2,\ C = c_1 c_2$), the gear-radius vector $\mathbf{r} = (a_1,b_1,c_1)\in\mathbb{Z}^3$ and turn-count vector $\mathbf{\omega} = (a_2,b_2,-c_2)\in\mathbb{Z}^3$ are strictly orthogonal ($\mathbf{r}\cdot\mathbf{\omega} = A+B-C = 0$) and together with the satellite normal $\mathbf{n} = \mathbf{r}\times\mathbf{\omega} = (-M_{BC},M_{CA},S_{AB})$ ($M_{BC} = b_1 c_2 + b_2 c_1$, $M_{CA} = c_1 a_2 + c_2 a_1$, $S_{AB} = a_1 b_2 - a_2 b_1$, $M_{AB} = a_1 b_2 + a_2 b_1$) form an orthogonal $3\times 3$ frame $\mathcal{F}_3$ with $\det\mathcal{F}_3 = M_{BC}^2 + M_{CA}^2 + S_{AB}^2 \equiv (a_1^2+b_1^2+c_1^2)(a_2^2+b_2^2+c_2^2) > 0$ and planar gear-slip law $C^2 + S_{AB}^2 = (A-B)^2 + M_{AB}^2 \equiv (a_1^2+b_1^2)(a_2^2+b_2^2)$ (where $S_{AB}=0$ for $\gcd(A,B)=1$ occurs solely at the Pell square degeneracy $(m,n)=(2,2)$; on $3+5^3=2^7=128$, $\|\mathbf{r}\|_2^2 = 1^2+5^2+4^2 = 42 = S_2(3)$, $S_{AB}=10$, $M_{AB}=40=k_{\mathrm{IV}}$, and $M_{CA}=44=C_{\mathrm{III}}$). On the cubic floor $k=3$ ($S_m = \widetilde{A}X^m + \widetilde{B}Y^m - \widetilde{C}Z^m$, $S_3 = 0$), the $3\times 3$ Hankel moment matrix $\mathcal{H}_3 = \begin{pmatrix} S_0 & S_1 & S_2 \\ S_1 & S_2 & 0 \\ S_2 & 0 & S_4 \end{pmatrix}$ obeys $-\det\mathcal{H}_3 = S_2^3 - S_4(S_0 S_2 - S_1^2) \equiv \widetilde{A}\widetilde{B}\widetilde{C}(X-Y)^2(Y-Z)^2(Z-X)^2 > 0$ (with $S_1 = S_3 = 0$ and $-\det\mathcal{H}_3 = 6\times 12^2 = 864$ on $3+5^3=2^7=128$) and the three Lagrange--Vandermonde bone-gap projectors $\widetilde{A}X(X-Y)(X-Z) \equiv S_1 YZ - S_2(Y+Z)$, $\widetilde{B}Y(Y-X)(Y-Z) \equiv S_1 ZX - S_2(Z+X)$, and $\widetilde{C}Z(Z-X)(Z-Y) \equiv S_2(X+Y) - S_1 XY$ ($\max(|S_1|,|S_2|)\ge 1$ for distinct bones $X,Y,Z$). Finally, coprimality $\gcd(A,B,C)=1$ defines a unique conjugate B\'ezout partner triple $u+v=w$ ($0\le u < A,\ 0 < v \le B,\ 0 < w < C$) forming $\mathcal{B}_{2\times 3} = \begin{pmatrix} A & B & C \\ u & v & w \end{pmatrix}$ with all three $2\times 2$ minors equal to $+1$ ($Av-Bu = Aw-Cu = Cv-Bw = +1$; yielding $(u,v,w) = (2,11,13)$ on $5+27=32$ and $(1,42,43)$ on $3+125=128$ with $v=42=\|\mathbf{r}\|_2^2$), while the bone-B\'ezout parent $(w_3,v_3)$ of $(Z_3,Y_3)$ ($Z_3 v_3 - Y_3 w_3 = +1$, $A' = \widetilde{C}_3 w_3^3 - \widetilde{B}_3 v_3^3$) forms the cube matrix $\mathcal{K}_{\mathrm{B\acute{e}zout}} = \begin{pmatrix} w_3^3 & -Z_3^3 \\ v_3^3 & -Y_3^3 \end{pmatrix}$ with Gauss-shell determinant $\det\mathcal{K}_{\mathrm{B\acute{e}zout}} = (Y_3 w_3+1)^3 - (Y_3 w_3)^3 = \Phi_3(Z_3 v_3, Y_3 w_3) > 0$ reconstructing the radical-bounded weights via $\mathcal{K}_{\mathrm{B\acute{e}zout}}(A,A')^T = \Phi_3(Z_3 v_3, Y_3 w_3)(\widetilde{B}_3,\widetilde{C}_3)^T$ ($\det\mathcal{K}_{\mathrm{B\acute{e}zout}} = 4^3-3^3 = 37$ on Lock~II, and level-$1$ collision invariant $I_1 = (C+2B)-A Y_3 w_3 = 333 = 9\times 37$ on Lock~IV).
\end{enumerate}
\end{theorem}

\section{Formal Verification in Lean~4}
\label{sec:lean4}

The complete self-contained formalization ($1\,670$ lines, zero external library dependencies beyond the standard Lean~4 core) resides in \texttt{RotatorCalculusPaper.lean} in the standalone open repository \url{https://github.com/KarataevDmitry/cuboctahedral-rotator-calculus} (licensed under strong reciprocal Copyleft \texttt{CC-BY-SA-4.0 / GPL-3.0-or-later}) and builds via \texttt{lake build}. In addition to the eight operator verification theorems (\texttt{thm1\_constant\_wronskian\_weierstrass\_lift\_and\_newton\_leibniz} through \texttt{thm8\_unified\_3x3\_pillai\_abc\_skeletal\_matrix\_and\_radical\_floor\_embedding}), the module proves six universally $\forall$-quantified algebraic and Diophantine-bound theorems across all integers $\mathbb{Z}$ (\texttt{thm9\_universal\_pillai\_abc\_differential\_bezout\_and\_pade\_laws} through \texttt{thm14\_master\_pillai\_and\_abc\_bilateral\_convergence\_chains}, which chains the cubic differential premise $\widetilde{A}X^3+\widetilde{B}Y^3-\widetilde{C}Z^3=0$, the bone-B\'ezout coprimality relation $Yw-Zv=1$, and the Mason--Stothers--B\'ezout coprimality relation $\widetilde{A}Xv_0-\widetilde{B}Yu_0=1$ via direct implication $\to$ into $A^*+B^*=C^*$, $W^* = X^2 Y^2 Z^2$, $-\det\mathcal{W}_{\mathrm{MS}} = X^4+Y^4+Z^4$, the bilateral Cramer convergence equations $w X^2 + \widetilde{B}\gamma = D_{\mathrm{clash}} Z$ and $v X^2 + \widetilde{C}\beta = D_{\mathrm{clash}} Y$, and the bilateral radical/gap bounds $C I_1 \le \widetilde{B}\widetilde{C}(Yw+Zv)$ and $Z^3(-I_1) \le A w^3(2Zv-Yw)$). Axiom audit (\texttt{\#print axioms}) of all fourteen theorems confirms zero non-constructive axioms (\texttt{0 sorryAx}, \texttt{0 Classical.choice}).

\section*{Acknowledgments and Declaration}
The author acknowledges formal verification partner \emph{Luch} for collaboration on the formal operator proofs. Large language model tools were used for assistance with \LaTeX{} formatting, English copy-editing, and translating mathematical proofs into Lean~4 tactic syntax; all mathematical definitions, theorem formulations, and proofs are verified by the Lean~4 formal proof kernel and are the full responsibility of the author.

\begin{thebibliography}{9}
\bibitem{Krause1986}
E.\,F.~Krause, \emph{Taxicab Geometry: An Adventure in Non-Euclidean Geometry}, Dover Publications, New York, 1986.

\bibitem{Thompson2011}
K.\,P.~Thompson, ``The nature of length, area, and volume in taxicab geometry,'' \emph{arXiv preprint} \texttt{arXiv:1101.2922} [math.MG], 2011.

\bibitem{Thompson2012}
K.\,P.~Thompson, ``Taxicab calculus: Trig derivatives,'' \emph{arXiv preprint} \texttt{arXiv:1205.3814v2} [math.MG], 2012.

\bibitem{Stroud1971}
A.\,H.~Stroud, \emph{Approximate Calculation of Multiple Integrals}, Prentice-Hall, Englewood Cliffs, NJ, 1971.

\bibitem{Moura2021}
L.~de~Moura and S.~Ullrich, ``The Lean~4 theorem prover and programming language,'' in \emph{Automated Deduction -- CADE~28}, Lecture Notes in Computer Science, vol.~12699, Springer, 2021, pp.~625--635.
\end{thebibliography}

\end{document}
